% ─────────────────────────────────────────────────────────────────────────────
% teoria_invariantes_klein.tex
% Overleaf: Menu -> Compiler -> pdfLaTeX
% ─────────────────────────────────────────────────────────────────────────────
\documentclass[11pt, a4paper]{article}

% -- Codificacion --------------------------------------------------------------
\usepackage[T1]{fontenc}
\usepackage[utf8]{inputenc}

% -- Fuentes (disponibles en TeX Live / Overleaf) -------------------------------
\usepackage[default, scale=0.95]{sourcesanspro}
\usepackage[scaled=0.82]{inconsolata}

% -- Geometria -------------------------------------------------------------------
\usepackage[
  a4paper,
  top        = 2.8cm,
  bottom     = 2.4cm,
  left       = 2.2cm,
  right      = 2.2cm,
  headheight = 30pt,
  headsep    = 14pt
]{geometry}

% -- Colores ----------------------------------------------------------------------
\usepackage[table]{xcolor}
\definecolor{cNavy}   {HTML}{0D1B3E}
\definecolor{cBlue}   {HTML}{2A4F9B}
\definecolor{cAccent} {HTML}{4A7FD4}
\definecolor{cMuted}  {HTML}{6B7FA8}
\definecolor{cCodeFG} {HTML}{1B3A70}
\definecolor{cCodeBG} {HTML}{EEF2FA}
\definecolor{cTHbg}   {HTML}{0D1B3E}
\definecolor{cTRalt}  {HTML}{F2F5FC}
\definecolor{cRule}   {HTML}{C0CADF}
\definecolor{cQuoteBG}{HTML}{F0F5FF}
\definecolor{cGood}   {HTML}{1E7A4C}
\definecolor{cBad}    {HTML}{B23A3A}

% -- Matematicas --------------------------------------------------------------------
\usepackage{amsmath}
\usepackage{amssymb}
\usepackage{mathtools}

% -- Tipografia y espaciado -----------------------------------------------------------
\usepackage{microtype}
\usepackage{parskip}
\setlength{\parindent}{0pt}
\setlength{\parskip}{5pt plus 1pt}

% -- Tablas --------------------------------------------------------------------------
\usepackage{booktabs}
\usepackage{tabularx}
\usepackage{colortbl}
\usepackage{array}
\usepackage{adjustbox}
\renewcommand{\arraystretch}{1.4}
\newcolumntype{L}[1]{>{\raggedright\arraybackslash}p{#1}}
\newcolumntype{C}[1]{>{\centering\arraybackslash}p{#1}}
\newcolumntype{R}[1]{>{\raggedleft\arraybackslash}p{#1}}

% -- Cajas de codigo y citas -----------------------------------------------------------
\usepackage{tcolorbox}
\tcbuselibrary{skins, breakable}

\tcbset{
  codebox/.style = {
    enhanced,
    colback    = cCodeBG,
    colframe   = cAccent,
    arc        = 3pt,
    leftrule   = 3pt,  rightrule  = 0.5pt,
    toprule    = 0.5pt, bottomrule = 0.5pt,
    left = 10pt, right = 10pt, top = 7pt, bottom = 7pt,
    fontupper  = \small\ttfamily\color{cCodeFG},
    breakable
  },
  quotebox/.style = {
    enhanced,
    colback    = cQuoteBG,
    colframe   = cAccent,
    arc        = 2pt,
    leftrule   = 3pt,  rightrule  = 0pt,
    toprule    = 0pt,  bottomrule = 0pt,
    left = 12pt, right = 10pt, top = 7pt, bottom = 7pt,
    fontupper  = \itshape\color{cCodeFG},
    breakable
  },
  defbox/.style = {
    enhanced,
    colback    = white,
    colframe   = cBlue,
    arc        = 2pt,
    leftrule   = 3pt,  rightrule  = 0.5pt,
    toprule    = 0.5pt, bottomrule = 0.5pt,
    left = 12pt, right = 10pt, top = 7pt, bottom = 7pt,
    breakable
  }
}

% -- Encabezados de seccion -------------------------------------------------------------
\usepackage{titlesec}

\titleformat{\section}
  {\color{cNavy}\Large\bfseries}
  {}{0em}{}
  [\vspace{2pt}{\color{cAccent}\hrule height 1.6pt}\vspace{4pt}]

\titleformat{\subsection}
  {\color{cBlue}\large\bfseries}
  {}{0em}{}
  [\vspace{1pt}{\color{cRule}\hrule height 0.5pt}\vspace{2pt}]

\titleformat{\subsubsection}
  {\color{cBlue}\normalsize\bfseries}
  {}{0em}{}

\titlespacing*{\section}      {0pt}{20pt}{8pt}
\titlespacing*{\subsection}   {0pt}{14pt}{4pt}
\titlespacing*{\subsubsection}{0pt}{10pt}{3pt}

% -- Cabecera y pie -----------------------------------------------------------------
\usepackage{fancyhdr}
\usepackage{tikz}
\usetikzlibrary{calc}

\pagestyle{fancy}
\fancyhf{}
\renewcommand{\headrulewidth}{0pt}

\fancyhead[C]{%
  \begin{tikzpicture}[remember picture, overlay]
    \fill[cNavy]
      (current page.north west)
      rectangle ($(current page.north east)+(0,-1.15cm)$);
    \fill[cAccent]
      (current page.north west)
      rectangle ($(current page.north west)+(0.28\paperwidth,-1.15cm)$);
    \node[white, anchor=west, font=\small\itshape] at
      ($(current page.north west)+(0.28\paperwidth+8pt,-0.58cm)$)
      {Go Ising \enspace\textbullet\enspace
       Experimento~08 \enspace\textbullet\enspace
       Teor\'ia de invariantes};
  \end{tikzpicture}%
}

\fancyfoot[L]{\small\color{cMuted}
  Leonardo Jim\'enez Mart\'inez \enspace\textbullet\enspace UNAM \enspace\textbullet\enspace 2026}
\fancyfoot[R]{\small\color{cMuted} P\'agina \thepage}
\def\footrule{{\color{cRule}\vspace{-3pt}\hrule width\headwidth height 0.4pt\vspace{2pt}}}

% -- Listas ----------------------------------------------------------------------------
\usepackage{enumitem}
\setlist[itemize]{
  leftmargin = 1.5em, itemsep = 2pt, topsep = 4pt,
  label = {\color{cAccent}\small$\bullet$}
}
\setlist[enumerate]{leftmargin=1.8em, itemsep=2pt, topsep=4pt}

% -- PDF metadata ------------------------------------------------------------------------
\usepackage[
  colorlinks = true,
  linkcolor  = cAccent,
  urlcolor   = cAccent,
  pdfauthor  = {Leonardo Jim\'enez Mart\'inez},
  pdftitle   = {Por qu\'e un grupo de Klein: teor\'ia de invariantes para el campo de Ising en Go},
  pdfsubject = {Go Ising -- Experimento 08 -- Teor\'ia de invariantes}
]{hyperref}

% -- Macros -------------------------------------------------------------------------------
\newcommand{\TH}[1]{\cellcolor{cTHbg}\color{white}\textbf{#1}}
\newcommand{\good}[1]{\color{cGood}\textbf{#1}}
\newcommand{\bad}[1]{\color{cBad}\textbf{#1}}

% =================================================================================
\begin{document}
\thispagestyle{fancy}

% -- Portada ----------------------------------------------------------------------------
\vspace*{4pt}
{\color{cNavy}\Huge\bfseries El grupo de Klein detr\'as}\\[3pt]
{\color{cNavy}\Large\bfseries de la degeneraci\'on algebraica}

\vspace{6pt}
{\color{cMuted}\large\itshape
  Teor\'ia de invariantes --- por qu\'e el campo de relajaci\'on solo ve
  4 de los 7 par\'ametros de \texttt{cubic\_mixed}}

\vspace{4pt}
{\color{cMuted}\small\itshape Generado autom\'aticamente \textbullet\ 2026-07-18}\\[1pt]
{\color{cMuted}\small\itshape
  Experimento 08 --- deriva formalmente un hallazgo del Experimento 07}

\vspace{8pt}
{\color{cAccent}\hrule height 2pt}
\vspace{10pt}

% =================================================================================
\section{Motivaci\'on: de un hallazgo ad hoc a una herramienta sistem\'atica}

En el informe de avance (Secci\'on~5.3) encontramos, revisando el
c\'odigo de \texttt{relaxation\_field}, que el campo que predice moyos
depende \'unicamente de 4 combinaciones de los 7 coeficientes de
\texttt{cubic\_mixed}:
\[
\Sigma a = a_1+a_2, \qquad \Sigma b = b_{11}+b_{22}, \qquad b_{12},
\qquad \Sigma c = c_{112}+c_{122}.
\]
Ese hallazgo se hizo \textbf{a mano}: expandimos $H(s,q)+H(q,s)$
algebraicamente y vimos qu\'e sobreviv\'ia. Funciona, pero no explica
\emph{por qu\'e} son exactamente esas 4 combinaciones y no otras, ni
qu\'e pasar\'ia con una plantilla distinta. Este informe deriva lo
mismo de manera \textbf{sistem\'atica y en un solo sentido}: primero
identificamos la \'unica simetr\'ia que hace falta para responder
"por qu\'e 4 y no 7" (posici\'on, $\sigma$), la resolvemos por
completo, y \textbf{despu\'es} a\~nadimos una segunda simetr\'ia
(color, $\tau$) para refinar la clasificaci\'on interna de esas 4 ---
sin volver a tocar el conteo principal.

% =================================================================================
\section{Los 7 polinomios de partida, uno por uno}
\label{sec:los7}

La plantilla \texttt{cubic\_mixed} es
\[
H(x,y) \;=\; a_1x + a_2y + b_{11}x^2 + b_{12}xy + b_{22}y^2
+ c_{112}x^2y + c_{122}xy^2,
\]
7 monomios, cada uno con su propio coeficiente libre:

\medskip
\begin{center}
\rowcolors{2}{white}{cTRalt}
\begin{tabular}{clc}
  \TH{\#} & \TH{Monomio} & \TH{Coeficiente} \\[1pt]
  1 & $x$     & $a_1$ \\
  2 & $y$     & $a_2$ \\
  3 & $x^2$   & $b_{11}$ \\
  4 & $xy$    & $b_{12}$ \\
  5 & $y^2$   & $b_{22}$ \\
  6 & $x^2y$  & $c_{112}$ \\
  7 & $xy^2$  & $c_{122}$ \\
\end{tabular}
\end{center}
\medskip

\'Estos son los 7 par\'ametros libres de cualquier Hamiltoniano de
esta familia --- el punto de partida completo. Las secciones que
siguen resuelven, en un solo camino hacia adelante, qu\'e les pasa a
estos 7 al pasar por \texttt{relaxation\_field}.

% =================================================================================
\section{¿Qu\'e es un grupo, aqu\'i?}

Un \textbf{grupo} es, para lo que necesitamos, una colecci\'on de
transformaciones que se pueden componer entre s\'i, donde:
\begin{itemize}
  \item hay una transformaci\'on "no hacer nada" (la \textbf{identidad} $e$);
  \item toda transformaci\'on tiene una \textbf{inversa} (deshacerla);
  \item componer transformaciones es \textbf{asociativo}.
\end{itemize}

Ejemplo m\'as simple: $\{+1,-1\}$ con la multiplicaci\'on usual es un
grupo de 2 elementos --- el m\'as peque\~no no trivial que existe. Es
exactamente el tama\~no de grupo que necesitamos primero.

% =================================================================================
\section{Paso 1: la simetr\'ia de posici\'on $\sigma$, y por qu\'e basta ella sola}
\label{sec:sigma}

\subsection{De d\'onde sale $\sigma$}

$H(x,y)$ punt\'ua la interacci\'on de un par de intersecciones
vecinas. Pero \textbf{no hay nada f\'isico que distinga cu\'al de las
dos es "la primera" ($i$) y cu\'al "la segunda" ($j$)} --- es una
elecci\'on de notaci\'on, no una propiedad del tablero. Por eso, en
\texttt{relaxation\_field} (\texttt{features.py:136}), cada
actualizaci\'on de un punto \textbf{siempre} usa la suma
$H(s,q)+H(q,s)$: ambos \'ordenes, sumados. Esa suma es, literalmente,
promediar sobre la transformaci\'on
\[
\sigma: (x,y) \;\longmapsto\; (y,x).
\]
Aplicarla dos veces regresa al original ($\sigma(\sigma(x,y))=(x,y)$)
--- $\sigma$ es una \textbf{involuci\'on}, y $\{e,\sigma\}$ es un
grupo de 2 elementos, el m\'as simple posible.

\subsection{Qu\'e le hace $\sigma$ a cada uno de los 7 monomios}

Un grupo de 2 elementos actuando sobre un espacio vectorial real
siempre lo separa en exactamente dos partes: la que \textbf{no cambia}
al aplicar la transformaci\'on (invariante) y la que \textbf{cambia de
signo} (anti-invariante) --- no hay una tercera opci\'on posible para
un grupo de orden 2. Aplicado monomio por monomio:

\medskip
\begin{center}
\rowcolors{2}{white}{cTRalt}
\begin{tabular}{lccc}
  \TH{Par de monomios} & \TH{Combinaci\'on $\sigma$-invariante} & \TH{Combinaci\'on $\sigma$-anti} & \TH{Total} \\[1pt]
  $\{x,y\}$       & $x{+}y$ (1)             & $x{-}y$ (1)     & 2 \\
  $\{x^2,y^2\}$   & $x^2{+}y^2$ (1)         & $x^2{-}y^2$ (1) & 2 \\
  $\{xy\}$        & $xy$ (1, ya sim\'etrico: $xy=yx$) & --- (0) & 1 \\
  $\{x^2y,xy^2\}$ & $x^2y{+}xy^2$ (1)       & $x^2y{-}xy^2$ (1) & 2 \\
  \hline
  \TH{Total} & \good{4} & \bad{3} & 7 \\
\end{tabular}
\end{center}
\medskip

El t\'ermino $xy$ es el \'unico sin pareja: como $xy=yx$ ya es
sim\'etrico por s\'i mismo, no genera una combinaci\'on
anti-invariante --- por eso el conteo final es 4 contra 3, no 3.5 y 3.5.

\subsection{Por qu\'e \texttt{relaxation\_field} solo ve la mitad invariante}

$H(s,q)+H(q,s)$ es, por construcci\'on, el doble de la proyecci\'on
sobre la parte $\sigma$-invariante: si $A(x,y)$ es cualquier
combinaci\'on \textbf{anti}-invariante (cambia de signo bajo $\sigma$),
entonces $A(s,q)+A(q,s) = A(s,q)-A(s,q) = 0$ exactamente --- se cancela
sin importar el valor de sus coeficientes. Verificado simb\'olicamente
con \texttt{sympy} sobre las tres combinaciones anti-invariantes de
\texttt{cubic\_mixed} ($x{-}y$, $x^2{-}y^2$, $x^2y{-}xy^2$, con
coeficientes libres $\Delta a=a_1{-}a_2$, $\Delta b=b_{11}{-}b_{22}$,
$\Delta c=c_{112}{-}c_{122}$): la suma da id\'enticamente cero.

\begin{tcolorbox}[quotebox]
  \textbf{Por qu\'e 4 y no 7, completo, con una sola simetr\'ia}: de
  los 7 monomios, 4 sobreviven al promedio $\sigma$ (combin\'andose en
  $\Sigma a=a_1{+}a_2$, $\Sigma b=b_{11}{+}b_{22}$, $b_{12}$ tal cual,
  y $\Sigma c=c_{112}{+}c_{122}$) y 3 se cancelan exactamente
  ($\Delta a$, $\Delta b$, $\Delta c$). Esta es la respuesta completa a
  "por qu\'e 4 y no 7" --- \textbf{no hace falta nada m\'as} que esta
  \'unica simetr\'ia de 2 elementos para derivarla.
\end{tcolorbox}

\subsection{Verificaci\'on emp\'irica: una predicci\'on falsable}

La teor\'ia predice algo m\'as fuerte que "solo importan 4
combinaciones": predice que un Hamiltoniano construido
\textbf{\'unicamente} con la parte anti-invariante ($\Delta a$,
$\Delta b$, o $\Delta c$ puros) deber\'ia dar $\Delta R^2=0$
\textbf{exacto}, sin importar la magnitud de sus coeficientes. Se
prob\'o sobre el cach\'e real de 20 partidas
(\texttt{cache\_full20\_cats}, radio~9):

\medskip
\begin{center}
\rowcolors{2}{white}{cTRalt}
\begin{tabular}{lr}
  \TH{Hamiltoniano} & \TH{$\Delta R^2$} \\[1pt]
  $\Delta a$ puro ($a_1{=}1,a_2{=}{-}1$, resto 0)         & \bad{0.000000} \\
  $\Delta b$ puro ($b_{11}{=}1,b_{22}{=}{-}1$, resto 0)   & \bad{0.000000} \\
  $\Delta c$ puro ($c_{112}{=}1,c_{122}{=}{-}1$, resto 0) & \bad{0.000000} \\
  Los tres combinados                                      & \bad{0.000000} \\
  \good{Control}: $\Sigma a$ puro ($a_1{=}1,a_2{=}1$)      & \good{0.000241} ($p=0.58$, distinto de 0) \\
\end{tabular}
\end{center}
\medskip

\begin{tcolorbox}[quotebox]
  Cero exacto, con precisi\'on de m\'aquina, para los cuatro
  Hamiltonianos puramente anti-invariantes --- no aproximadamente
  cero, \textbf{cero}. El control ($\Sigma a$ puro, que s\'i debe
  sobrevivir) da un valor peque\~no pero distinto de cero. La
  predicci\'on te\'orica se confirma sin ambig\"uedad sobre datos
  reales, y con esto qued\'o completamente resuelta la pregunta
  central del informe.
\end{tcolorbox}

% =================================================================================
\section{Paso 2: una segunda simetr\'ia, $\tau$, para refinar --- no para recontar}
\label{sec:tau}

Hasta aqu\'i, todo se deriv\'o con \textbf{una sola} simetr\'ia
($\sigma$). Esta secci\'on a\~nade una \textbf{segunda}, independiente,
para responder una pregunta distinta: de las 4 combinaciones que
sobrevivieron ($\Sigma a,\Sigma b,b_{12},\Sigma c$), ¿hay alguna
subestructura interna? ¿Por qu\'e $H_{M1}$ usa solo $\Sigma a,\Sigma c$
y Alvarado solo $b_{12}$?

\subsection{La simetr\'ia de color, $\tau$}

Tampoco hay nada f\'isico que distinga "negro" de "blanco" como
etiquetas --- son convenciones ($s=-1$ vs.\ $s=+1$). Invertir todos los
colores del tablero corresponde a
\[
\tau: (x,y) \;\longmapsto\; (-x,-y).
\]
Al igual que $\sigma$, es una involuci\'on ($\tau^2=e$).

\subsection{$\sigma$ y $\tau$ conmutan}

\[
\sigma(\tau(x,y)) = \sigma(-x,-y) = (-y,-x)
\qquad = \qquad
\tau(\sigma(x,y)) = \tau(y,x) = (-y,-x).
\]
Verificado tambi\'en sobre $H$ completo con \texttt{sympy}: aplicar
$\sigma$ despu\'es de $\tau$ da t\'ermino a t\'ermino lo mismo que
$\tau$ despu\'es de $\sigma$.

\subsection{Por qu\'e $\{e,\sigma,\tau,\sigma\tau\}$ es el grupo de Klein, no $\mathbb{Z}_4$}
\label{sec:klein}

Juntas, $\sigma$ y $\tau$ generan un grupo de 4 elementos,
$G=\{e,\sigma,\tau,\sigma\tau\}$. Hay \textbf{dos} grupos distintos de
4 elementos (salvo isomorfismo), y hay que distinguir cu\'al es \'este:

\begin{itemize}
  \item $\mathbb{Z}_4$ (c\'iclico): tiene un elemento de \textbf{orden 4}.
  \item $\mathbb{Z}_2\times\mathbb{Z}_2$ (\textbf{grupo de Klein}, $V_4$):
        \textbf{todos} sus elementos no triviales tienen orden 2.
\end{itemize}

Aqu\'i, $\sigma^2=e$, $\tau^2=e$, y $(\sigma\tau)^2=\sigma^2\tau^2=e$
(usando que conmutan) --- \textbf{los tres elementos no triviales
tienen orden 2}: es el grupo de Klein. Tabla de multiplicaci\'on
completa:

\medskip
\begin{center}
\rowcolors{2}{white}{cTRalt}
\begin{tabular}{c|cccc}
  $\cdot$ & $e$ & $\sigma$ & $\tau$ & $\sigma\tau$ \\
  \hline
  $e$          & $e$          & $\sigma$      & $\tau$        & $\sigma\tau$ \\
  $\sigma$     & $\sigma$     & $e$           & $\sigma\tau$  & $\tau$ \\
  $\tau$       & $\tau$       & $\sigma\tau$  & $e$           & $\sigma$ \\
  $\sigma\tau$ & $\sigma\tau$ & $\tau$        & $\sigma$      & $e$ \\
\end{tabular}
\end{center}
\medskip

Cada elemento aparece exactamente una vez por fila/columna (grupo
v\'alido) y la diagonal es puro $e$ --- la firma exacta del grupo de
Klein.

\begin{tcolorbox}[quotebox]
  \textbf{Por qu\'e importa que sea Klein y no c\'iclico}: en
  $\mathbb{Z}_4$ las representaciones irreducibles usan n\'umeros
  complejos (ra\'ices de la unidad) porque hay un elemento de orden 4.
  En Klein, como todo tiene orden $\leq2$, todas las representaciones
  son \textbf{reales}, con valores $\pm1$ \'unicamente --- justo lo que
  necesitamos, porque nuestros coeficientes son n\'umeros reales.
\end{tcolorbox}

\subsection{Las 4 piezas de Klein, y d\'onde queda cada mon\'omio}

Como $G$ es abeliano de orden 4, tiene exactamente 4 representaciones
irreducibles, todas de dimensi\'on 1 --- cada una es un \textbf{car\'acter}
que asigna $\pm1$ a $\sigma$ y a $\tau$ por separado:

\medskip
\begin{center}
\rowcolors{2}{white}{cTRalt}
\begin{tabular}{lccl}
  \TH{Car\'acter} & \TH{$\sigma\mapsto$} & \TH{$\tau\mapsto$} & \TH{Significado} \\[1pt]
  $\chi_{++}$ & $+1$ & $+1$ & sim\'etrico en posici\'on y en color \\
  $\chi_{+-}$ & $+1$ & $-1$ & sim\'etrico en posici\'on, impar en color \\
  $\chi_{-+}$ & $-1$ & $+1$ & antisim\'etrico en posici\'on, par en color \\
  $\chi_{--}$ & $-1$ & $-1$ & antisim\'etrico en ambas \\
\end{tabular}
\end{center}
\medskip

El operador de Reynolds $P_\chi(H)=\frac{1}{4}\sum_{g\in G}\chi(g)\,(g{\cdot}H)$
proyecta sobre cada pieza. Aplic\'andolo a \texttt{cubic\_mixed}
(verificado simb\'olicamente, suma exacta $=H$):

\begin{tcolorbox}[codebox]
$P_{++}(H) = \dfrac{\Sigma b}{2}(x^2{+}y^2) + b_{12}\,xy$
\hfill \textit{(2 dim.)} \\[4pt]
$P_{+-}(H) = \dfrac{\Sigma a}{2}(x{+}y) + \dfrac{\Sigma c}{2}(x^2y{+}xy^2)$
\hfill \textit{(2 dim.)} \\[4pt]
$P_{-+}(H) = \dfrac{\Delta b}{2}(x^2{-}y^2)$
\hfill \textit{(1 dim.)} \\[4pt]
$P_{--}(H) = \dfrac{\Delta a}{2}(x{-}y) + \dfrac{\Delta c}{2}(x^2y{-}xy^2)$
\hfill \textit{(2 dim.)}
\end{tcolorbox}

N\'otese: $P_{++}\oplus P_{+-}$ (dimensi\'on $2+2=4$) es exactamente la
parte $\sigma$-invariante de la Secci\'on~\ref{sec:sigma} --- $\tau$ no
cambi\'o cu\'antas dimensiones sobreviven, solo las \textbf{parti\'o en
dos} seg\'un el color. Y $P_{-+}\oplus P_{--}$ ($1+2=3$) es la parte
$\sigma$-anti que ya sab\'iamos que muere, ahora tambi\'en partida en
dos por color (irrelevante para \texttt{relaxation\_field}, pero
consistente).

% =================================================================================
\section{Tabla final: los 7 monomios, de principio a fin}

Uniendo ambos pasos, as\'i queda cada uno de los 7 monomios
originales:

\medskip
\begin{center}
\rowcolors{2}{white}{cTRalt}
\begin{tabular}{clcccl}
  \TH{\#} & \TH{Monomio} & \TH{Coef.} & \TH{Combinaci\'on que sobrevive} & \TH{Pieza de Klein} & \TH{¿Sobrevive?} \\[1pt]
  1 & $x$    & $a_1$     & $\Sigma a = a_1{+}a_2$         & $P_{+-}$ & \good{S\'i} \\
  2 & $y$    & $a_2$     & $\Sigma a = a_1{+}a_2$         & $P_{+-}$ & \good{S\'i} \\
  3 & $x^2$  & $b_{11}$  & $\Sigma b = b_{11}{+}b_{22}$   & $P_{++}$ & \good{S\'i} \\
  4 & $xy$   & $b_{12}$  & $b_{12}$ (sin pareja, intacto) & $P_{++}$ & \good{S\'i} \\
  5 & $y^2$  & $b_{22}$  & $\Sigma b = b_{11}{+}b_{22}$   & $P_{++}$ & \good{S\'i} \\
  6 & $x^2y$ & $c_{112}$ & $\Sigma c = c_{112}{+}c_{122}$ & $P_{+-}$ & \good{S\'i} \\
  7 & $xy^2$ & $c_{122}$ & $\Sigma c = c_{112}{+}c_{122}$ & $P_{+-}$ & \good{S\'i} \\
\end{tabular}
\end{center}
\medskip

\begin{tcolorbox}[defbox]
\textbf{Las 4 combinaciones que sobreviven, tal cual:}
\begin{center}
\rowcolors{2}{white}{cTRalt}
\begin{tabular}{lll}
  \TH{Combinaci\'on} & \TH{F\'ormula} & \TH{Viene de} \\[1pt]
  $\Sigma a$ & $a_1+a_2$       & coeficientes de $x$ y $y$ \\
  $\Sigma b$ & $b_{11}+b_{22}$ & coeficientes de $x^2$ y $y^2$ \\
  $b_{12}$   & $b_{12}$        & coeficiente de $xy$ \\
  $\Sigma c$ & $c_{112}+c_{122}$ & coeficientes de $x^2y$ y $xy^2$ \\
\end{tabular}
\end{center}
\end{tcolorbox}

\begin{center}
\rowcolors{2}{white}{cTRalt}
\begin{tabular}{lll}
  \TH{Las 3 que mueren} & \TH{F\'ormula} & \TH{Viene de} \\[1pt]
  $\Delta a$ & $a_1-a_2$       & coeficientes de $x$ y $y$ \\
  $\Delta b$ & $b_{11}-b_{22}$ & coeficientes de $x^2$ y $y^2$ \\
  $\Delta c$ & $c_{112}-c_{122}$ & coeficientes de $x^2y$ y $xy^2$ \\
\end{tabular}
\end{center}

% =================================================================================
\section{Lo que $\tau$ explica que $\sigma$ sola no explicaba}
\label{sec:paraque2}

\subsection{Correcci\'on importante: $H_{M1}$ usa \textbf{dos} particiones, no una}

Una primera lectura sugiere que, como \texttt{sparse\_cubic}/$H_{M1}$
solo usa $a_1,a_2,c_{112},c_{122}$ (nunca $b_{11},b_{12},b_{22}$), su
plantilla vivir\'ia \'integramente en $P_{+-}$. \textbf{Eso es
incorrecto} si se toma $H_{M1}$ con sus coeficientes reales
--- hay que revisar las \textbf{dos} particiones $\tau$-impares
($P_{+-}$ \emph{y} $P_{--}$), no solo una. Descomponiendo
$H_{M1}=x+2y-x^2y-xy^2$ (verificado con \texttt{sympy}):

\begin{tcolorbox}[codebox]
$P_{+-}(H_{M1}) = \dfrac{3}{2}(x{+}y) \,-\, 1\cdot(x^2y{+}xy^2)$
\hfill ($\Sigma a{=}3,\ \Sigma c{=}{-}2$) \\[4pt]
$P_{--}(H_{M1}) = -\dfrac{1}{2}(x{-}y) \,+\, 0\cdot(x^2y{-}xy^2)$
\hfill ($\Delta a{=}{-}1,\ \Delta c{=}0$) \\[4pt]
$P_{++}(H_{M1}) = 0, \qquad P_{-+}(H_{M1}) = 0$
\end{tcolorbox}

$H_{M1} = P_{+-}(H_{M1}) + P_{--}(H_{M1})$, y \textbf{ninguna de las
dos partes es cero}: $a_1=1\neq a_2=2$, as\'i que $\Delta a=a_1-a_2=-1$
no se anula. $H_{M1}$ s\'i es completamente $\tau$-\textbf{impar}
(vive en la uni\'on $P_{+-}\oplus P_{--}$, nunca toca $P_{++}$ ni
$P_{-+}$ --- eso s\'i es correcto, y por eso $H(-x,-y)=-H(x,y)$ se
cumple exactamente, verificado). Pero \textbf{dentro} de ese lado
impar-en-color, $H_{M1}$ se reparte entre las dos particiones de
$\sigma$: la mayor parte en $P_{+-}$ (lo que \texttt{relaxation\_field}
s\'i ve), y una parte menor pero no nula en $P_{--}$ (lo que
\texttt{relaxation\_field} \textbf{nunca} ve --- se cancela en
$H(s,q)+H(q,s)$ sin importar que $H_{M1}$, como polinomio completo,
s\'i la tenga).

\begin{tcolorbox}[quotebox]
  \textbf{Esto no es un caso aislado}: se verific\'o que
  $H_{OPT\_A}$ (el \'optimo de Fase~A) \textbf{tambi\'en} tiene
  componente $P_{--}$ no nula ($\Delta a\approx-0.340$,
  $\Delta c\approx1.105$). Ning\'un Hamiltoniano probado hasta ahora en
  este experimento vive puramente en $P_{+-}$ --- la clasificaci\'on
  por "tipo" (Secci\'on siguiente) describe correctamente el signo de
  $\Sigma a,\Sigma c$ (lo que importa para $\Delta R^2$), pero
  \textbf{no} implica que el resto de $H$ sea cero; solo que es
  invisible para \texttt{relaxation\_field}.
\end{tcolorbox}

\subsection{¿Cu\'ando es puro y cu\'ando mixto? La regla general}
\label{sec:regla-mixto}

Lo anterior no es una peculiaridad de $H_{M1}$: es la regla para
\textbf{casi cualquier} Hamiltoniano. Cada par de coeficientes
originales ($a_1,a_2$), ($b_{11},b_{22}$), ($c_{112},c_{122}$) se
reparte en su suma ($\Sigma$, sobrevive) y su resta ($\Delta$, muere)
--- un Hamiltoniano es \textbf{puro} en una sola pieza de Klein
\'unicamente cuando todas sus $\Delta$ correspondientes son cero
exactamente:

\medskip
\begin{center}
\rowcolors{2}{white}{cTRalt}
\begin{tabular}{lll}
  \TH{Caso} & \TH{¿Por qu\'e es puro?} & \TH{Pieza} \\[1pt]
  Alvarado ($b_{12}=1$, resto 0) & $a_1{=}a_2{=}0$, $b_{11}{=}b_{22}{=}0$: & $P_{++}$ puro \\
   & todos los $\Delta$ son 0 trivialmente & \\
  Clase unitaria $(+1,+1,-1,-1)$ & $a_1{=}a_2{=}1\Rightarrow\Delta a{=}0$; & $P_{+-}$ puro \\
  (signos de $H_{M1}$, del barrido) & $c_{112}{=}c_{122}{=}{-}1\Rightarrow\Delta c{=}0$ & \\
  $H_{M1}$ real ($a_1{=}1,a_2{=}2$) & $a_1\neq a_2\Rightarrow\Delta a{=}{-}1\neq0$ & $P_{+-}\oplus P_{--}$, mixto \\
  $H_{OPT\_A}$ (optimizado) & $a_1\neq a_2$ y $c_{112}\neq c_{122}$ & $P_{+-}\oplus P_{--}$, mixto \\
  Cualquier candidato aleatorio & coeficientes continuos: la & casi seguro mixto \\
  (los 44 + 25 de este experimento) & probabilidad de $a_1{=}a_2$ exacto es 0 & \\
\end{tabular}
\end{center}
\medskip

\begin{tcolorbox}[quotebox]
  \textbf{Conclusi\'on pr\'actica}: la pureza (vivir en una sola pieza
  de Klein) es la \textbf{excepci\'on}, no la norma --- ocurre solo
  cuando los coeficientes se eligieron a prop\'osito con esa simetr\'ia
  extra (como Alvarado, o las clases unitarias del barrido con
  coeficientes iguales por pareja). $H_{M1}$, $H_{OPT\_A}$, y
  virtualmente todos los candidatos aleatorios de \texttt{cubic\_mixed}
  y \texttt{sparse\_cubic} son mixtos: tienen una parte que
  \texttt{relaxation\_field} ve (la proyecci\'on $\sigma$-invariante) y
  una parte que nunca ve, pero que sigue siendo parte real del
  polinomio.
\end{tcolorbox}

\subsubsection{Conteo exacto sobre los 47 Hamiltonianos catalogados}

Se calcul\'o $\Delta a=a_1{-}a_2$, $\Delta b=b_{11}{-}b_{22}$,
$\Delta c=c_{112}{-}c_{122}$ para \textbf{cada uno} de los 47
Hamiltonianos del cat\'alogo del visor (los 44 \texttt{cubic\_mixed} +
$H_{M1}$ + Alvarado + $H_{OPT\_A}$), no solo para los 3 ejemplos
citados arriba:

\medskip
\begin{center}
\rowcolors{2}{white}{cTRalt}
\begin{tabular}{lr}
  \TH{Categor\'ia} & \TH{Cantidad (de 47)} \\[1pt]
  \good{Puros} (los 3 $\Delta$ son exactamente 0) & \good{1} \\
  \bad{Mixtos} (al menos un $\Delta\neq0$) & \bad{46} \\
\end{tabular}
\end{center}
\medskip

\begin{tcolorbox}[quotebox]
  El \'unico Hamiltoniano puro del cat\'alogo completo es
  \textbf{Alvarado} ($H_{0001}$) --- y solo porque su plantilla
  \texttt{quadratic} fija $a_1=a_2=0$ y $b_{11}=b_{22}=0$
  trivialmente, no por dise\~no expl\'icito. Los otros 46 --- incluido
  $H_{M1}$ --- son mixtos: tienen una parte $\sigma$-invariante (la que
  $\Delta R^2$ mide) y una parte $\sigma$-anti que \texttt{relaxation\_field}
  nunca ve, pero que existe como parte real del polinomio. Ser
  "mixto" no es la excepci\'on entre los candidatos aleatorios --- es
  virtualmente universal.
\end{tcolorbox}

\textbf{Consecuencia para la clasificaci\'on por "tipo" del visor}: las
etiquetas ("tipo $H_{M1}$", "tipo espejo", \ldots) se calculan
\'unicamente a partir del signo de $\Sigma a,\Sigma b,b_{12},\Sigma c$
--- nunca miran $\Delta a,\Delta b,\Delta c$. Dos Hamiltonianos con el
mismo signo de $\Sigma a,\Sigma c$ pero distinto $\Delta a,\Delta c$
reciben la \textbf{misma} etiqueta de tipo, aunque sean polinomios
distintos como objetos matem\'aticos completos. La etiqueta describe
correctamente la parte relevante para $\Delta R^2$; no pretende
describir el Hamiltoniano completo.

\subsection{Qu\'e s\'i explica $\tau$, entonces}

$\tau$ divide el espacio en dos mitades por color (par/impar), y
$\sigma$ en dos mitades por posici\'on (sim\'etrico/anti) ---
\textbf{independientes entre s\'i}. $P_{+-}$ y $P_{--}$ juntas
($\Sigma a,\Sigma c,\Delta a,\Delta c$) son exactamente la mitad
$\tau$-impar --- la \'unica que puede usar \texttt{sparse\_cubic}, ya
que su plantilla no incluye t\'erminos $b$. Que $H_{M1}$ cargue
tambi\'en una parte $\Delta a$ no nula no rompe esa pertenencia a la
mitad impar; solo confirma que, \textbf{dentro} de esa mitad, $H_{M1}$
no est\'a alineado exactamente con el eje $\sigma$-invariante. $P_{++}$
(par en color) contiene $\Sigma b$ y $b_{12}$ --- Alvarado ($H=xy$)
vive aqu\'i, y en ese caso s\'i por completo, porque $b_{11}=b_{22}=0$
hace que $\Delta b=0$ trivialmente.

La clasificaci\'on por "tipo de polinomio" ya construida en el visor
(tipo $H_{M1}$, tipo espejo, cuadr\'atico puro, nulo) describe
correctamente el signo de $(\Sigma a,\Sigma b,b_{12},\Sigma c)$ ---
la parte que importa para $\Delta R^2$ --- pero, a la luz de esta
secci\'on, debe leerse como "la proyecci\'on $\sigma$-invariante tiene
este signo", no como "el Hamiltoniano completo es esto".

% =================================================================================
\section{Implicaci\'on pr\'actica y alcance}

\textbf{Para Fase~B}: buscar directamente en las 4 dimensiones
efectivas ($\Sigma a,\Sigma b,b_{12},\Sigma c$) en vez de las 7 crudas
es estrictamente m\'as eficiente, con garant\'ia te\'orica (no solo
emp\'irica) de que no se pierde nada.

\textbf{Generalidad de la herramienta}: el mismo procedimiento ---
identificar las simetr\'ias del problema, reconocer el grupo que
generan, descomponer en representaciones irreducibles, proyectar ---
aplica a cualquier plantilla de Hamiltoniano futura, no solo
\texttt{cubic\_mixed}.

\begin{tcolorbox}[quotebox]
  \textbf{Advertencia de alcance (igual que en el informe principal)}:
  esta reducci\'on aplica \'unicamente al pipeline de predicci\'on de
  moyo (\texttt{relaxation\_field}). El an\'alisis topol\'ogico de
  $\Gamma(H)$ (nodos $A_1$, puntos cr\'iticos) punt\'ua $H(x,y)$
  directamente, sin promediar sobre $\sigma$ --- ah\'i s\'i importan
  los 7 coeficientes por separado.
\end{tcolorbox}

\vfill
\noindent{\color{cRule}\hrule height 0.5pt}\vspace{4pt}
{\color{cMuted}\small\itshape
  Complementa:~\texttt{experiments/07\_moyo\_dataset/.../avance\_experimento07.tex} (Secci\'on~5.3)
  \quad\textbullet\quad
  Verificaci\'on simb\'olica:~\texttt{sympy}
  \quad\textbullet\quad
  Verificaci\'on emp\'irica:~\texttt{cache\_full20\_cats} (20 partidas, radio~9, generado en exp.~07)
}

\end{document}
